% GENERATED FILE. Edit canonical lessons, then run npm run generate:books.
% canonical-source-hash: fnv1a64:bb9a7a68152b0f0e
% canonical-lessons: RU-C182-teper, RU-C182-kanikuly, RU-C182-den-rozhdeniya, RU-C182-vykhodnoy, RU-C182-vovremya

\chapter{Nowadays, School holidays, Birthday, Day off, On time}
\label{ch:ru-nowadays-school-holidays-birthday-day-off-on-time}
\hlchaptermodality{182}

\begin{chapteropening}
\textbf{By the end of this chapter:} \emph{I can say nowadays, school holidays, a birthday, a day off and on time.}

Complete the last lesson of chapter 182: I can say nowadays, school holidays, a birthday, a day off and on time.
\end{chapteropening}

\section[\texorpdfstring{\ru{теперь} (tepér)}{теперь (tepér)}]{\ru{теперь} (tepér) — nowadays, now (after a change)}

\label{lesson:RU-C182-teper}



\begin{quote}
Before the new one: say the Russian for April, then the Russian for May.
\end{quote}

\subsection*{You'll want to know: \ru{теперь}}
\textbf{\ru{теперь}} (\emph{tepér}) — \textquotedblleft{}nowadays, now (after a change)\textquotedblright{}.

\begin{grammarlens}[title={What you've built}]
One more word for this chapter.
\end{grammarlens}

\subsection*{Your turn}
\begin{itemize}
\raggedright
  \item \emph{Say it:} \emph{tepér}
  \item \emph{Say it:} \emph{tepér}, once more
  \item \emph{Say it:} say \emph{tepér}
  \item \emph{Recall:} say the Russian for April, then the Russian for May, then say \emph{tepér} again
\end{itemize}

\begin{tcolorbox}[breakable,colback=teal!4,colframe=teal!35!black,title={Before you move on}]
What does \ru{теперь} mean? (\textquotedblleft{}Nowadays, now (after a change)\textquotedblright{}.) Say it once more.
\end{tcolorbox}
\section[\texorpdfstring{\ru{каникулы} (kaníkuly)}{каникулы (kaníkuly)}]{\ru{каникулы} (kaníkuly) — school holidays}

\label{lesson:RU-C182-kanikuly}



\begin{quote}
Before the new one: say the Russian for May, then the Russian for nowadays.
\end{quote}

\subsection*{You'll want to know: \ru{каникулы}}
\textbf{\ru{каникулы}} (\emph{kaníkuly}) — \textquotedblleft{}school holidays\textquotedblright{}.

\begin{grammarlens}[title={What you've built}]
One more word for this chapter.
\end{grammarlens}

\subsection*{Your turn}
\begin{itemize}
\raggedright
  \item \emph{Say it:} \emph{kaníkuly}
  \item \emph{Say it:} \emph{kaníkuly}, once more
  \item \emph{Say it:} say \emph{kaníkuly}
  \item \emph{Recall:} say the Russian for May, then the Russian for nowadays, then say \emph{kaníkuly} again
\end{itemize}

\begin{tcolorbox}[breakable,colback=teal!4,colframe=teal!35!black,title={Before you move on}]
What does \ru{каникулы} mean? (\textquotedblleft{}School holidays\textquotedblright{}.) Say it once more.
\end{tcolorbox}
\section[\texorpdfstring{\ru{день} \ru{рождения} (den rozhdéniya)}{день рождения (den rozhdéniya)}]{\ru{день} \ru{рождения} (den rozhdéniya) — a birthday}

\label{lesson:RU-C182-den-rozhdeniya}



\begin{quote}
Before the new one: say the Russian for nowadays, then the Russian for school holidays.
\end{quote}

\subsection*{You'll want to know: \ru{день} \ru{рождения}}
\textbf{\ru{день} \ru{рождения}} (\emph{den rozhdéniya}) — \textquotedblleft{}a birthday\textquotedblright{}.

\begin{grammarlens}[title={What you've built}]
One more word for this chapter.
\end{grammarlens}

\subsection*{Your turn}
\begin{itemize}
\raggedright
  \item \emph{Say it:} \emph{den rozhdéniya}
  \item \emph{Say it:} \emph{den rozhdéniya}, once more
  \item \emph{Say it:} say \emph{den rozhdéniya}
  \item \emph{Recall:} say the Russian for nowadays, then the Russian for school holidays, then say \emph{den rozhdéniya} again
\end{itemize}

\begin{tcolorbox}[breakable,colback=teal!4,colframe=teal!35!black,title={Before you move on}]
What does \ru{день} \ru{рождения} mean? (\textquotedblleft{}A birthday\textquotedblright{}.) Say it once more.
\end{tcolorbox}
\section[\texorpdfstring{\ru{выходной} (vykhodnóy)}{выходной (vykhodnóy)}]{\ru{выходной} (vykhodnóy) — a day off}

\label{lesson:RU-C182-vykhodnoy}



\begin{quote}
Before the new one: say the Russian for school holidays, then the Russian for a birthday.
\end{quote}

\subsection*{You'll want to know: \ru{выходной}}
\textbf{\ru{выходной}} (\emph{vykhodnóy}) — \textquotedblleft{}a day off\textquotedblright{}.

\begin{grammarlens}[title={What you've built}]
One more word for this chapter.
\end{grammarlens}

\subsection*{Your turn}
\begin{itemize}
\raggedright
  \item \emph{Say it:} \emph{vykhodnóy}
  \item \emph{Say it:} \emph{vykhodnóy}, once more
  \item \emph{Say it:} say \emph{vykhodnóy}
  \item \emph{Recall:} say the Russian for school holidays, then the Russian for a birthday, then say \emph{vykhodnóy} again
\end{itemize}

\begin{tcolorbox}[breakable,colback=teal!4,colframe=teal!35!black,title={Before you move on}]
What does \ru{выходной} mean? (\textquotedblleft{}A day off\textquotedblright{}.) Say it once more.
\end{tcolorbox}
\section[\texorpdfstring{\ru{вовремя} (vóvremya)}{вовремя (vóvremya)}]{\ru{вовремя} (vóvremya) — on time}

\label{lesson:RU-C182-vovremya}



\begin{quote}
Before the new one: say the Russian for a birthday, then the Russian for a day off.
\end{quote}

\subsection*{You'll want to know: \ru{вовремя}}
\textbf{\ru{вовремя}} (\emph{vóvremya}) — \textquotedblleft{}on time\textquotedblright{}.

\begin{grammarlens}[title={What you've built}]
One more word for this chapter.
\end{grammarlens}

\subsection*{Your turn}
\begin{itemize}
\raggedright
  \item \emph{Say it:} \emph{vóvremya}
  \item \emph{Say it:} \emph{vóvremya}, once more
  \item \emph{Say it:} say \emph{vóvremya}
  \item \emph{Recall:} say the Russian for a birthday, then the Russian for a day off, then say \emph{vóvremya} again
\end{itemize}

\begin{tcolorbox}[breakable,colback=teal!4,colframe=teal!35!black,title={Before you move on}]
What does \ru{вовремя} mean? (\textquotedblleft{}On time\textquotedblright{}.) Say it once more.
\end{tcolorbox}
